\documentclass[tikz,border=6pt]{standalone}

\usepackage[T1]{fontenc}
\usepackage{lmodern}
\usepackage{xcolor}
\usepackage{amssymb} % \checkmark

\usepackage{tikz}
\usetikzlibrary{arrows.meta,calc,decorations.pathreplacing,positioning}

% NOTE:
% This is a standalone figure intended to mirror the style of the legacy
% RobotClient diagram, but reflect the architecture in:
%   src/lerobot/async_inference/robot_client_alternative.py

\begin{document}
\begin{tikzpicture}[
  font=\small,
  laneLine/.style={line width=0.5pt},
  axisArrow/.style={-{Stealth[length=2.0mm]}, line width=0.6pt},
  popArrow/.style={-{Stealth[length=1.8mm]}, line width=0.5pt},
  chunkBox/.style={draw=black, line width=0.6pt, rounded corners=2pt, fill=white},
  overlapFill/.style={fill=blue!14, draw=none, opacity=0.60},
  labelText/.style={align=left},
  braceLabel/.style={midway, fill=white, inner sep=1pt},
]

% ----------------------------
% Layout constants
% ----------------------------
\def\xRecv{0.0}
\def\xCtrl{3.8}
\def\xAxisStart{6.4}
\def\xAxisEnd{16.6}

\def\yTop{0.0}
\def\yBottom{-11.4}

% Chunk geometry
\def\slotW{0.55}
\def\slotH{0.55}

% ----------------------------
% Title
% ----------------------------
\node[font=\large] at ($(0,0)+(8.3,0.75)$) {RobotClientAlternative};

% ----------------------------
% Swimlanes (left)
% ----------------------------
\draw[laneLine] (\xRecv,\yTop) -- (\xRecv,\yBottom);
\draw[laneLine] (\xCtrl,\yTop) -- (\xCtrl,\yBottom);

\node[font=\normalsize] at (\xRecv,0.25) {Action receiving};
\node[font=\normalsize] at (\xRecv,-0.15) {thread};
\node[font=\normalsize] at (\xCtrl,0.25) {Control loop};
\node[font=\normalsize] at (\xCtrl,-0.15) {thread};

\node at (\xRecv,\yBottom-0.55) {Timestamp, $t$};
\node at (\xCtrl,\yBottom-0.55) {Timestamp, $t$};

% ----------------------------
% Timeline annotations (left)
% ----------------------------
% Receiver events
\node[labelText, anchor=west] at (\xRecv+0.25,-1.4) {First action chunk received\\(enqueued into incoming queue)};
\node[labelText, anchor=west] at (\xRecv+0.25,-7.2) {New action chunk received\\(enqueued into incoming queue)};

% Control-loop events (key differences vs legacy)
\node[labelText, anchor=west] at (\xCtrl+0.25,-1.0) {Drain incoming chunks};
\node[labelText, anchor=west] at (\xCtrl+0.25,-1.9) {Action performed};
\node[labelText, anchor=west] at (\xCtrl+0.25,-2.8) {Observation not sent};
\node[labelText, anchor=west] at (\xCtrl+0.25,-3.8) {Action performed};
\node[labelText, anchor=west] at (\xCtrl+0.25,-4.7) {Observation sent};
\node[labelText, anchor=west] at (\xCtrl+0.25,-5.7) {Action performed};
\node[labelText, anchor=west] at (\xCtrl+0.25,-6.6) {Observation sent};

\node[labelText, anchor=west] at (\xCtrl+0.25,-7.9) {Drain incoming chunks};
\node[labelText, anchor=west] at (\xCtrl+0.25,-8.8) {Action performed};
\node[labelText, anchor=west] at (\xCtrl+0.25,-9.7) {Observation not sent};

% Markers (sent/not-sent)
\newcommand{\MarkX}[2]{\node at (#1,#2) {\textcolor{red}{\Large $\times$}};}
\newcommand{\MarkOK}[2]{\node at (#1,#2) {\textcolor{green!60!black}{\Large \checkmark}};}
\MarkX{\xCtrl+2.9}{-2.8}
\MarkOK{\xCtrl+2.9}{-4.7}
\MarkOK{\xCtrl+2.9}{-6.6}
\MarkX{\xCtrl+2.9}{-9.7}

% ----------------------------
% Right side: timestep axis
% ----------------------------
\draw[axisArrow] (\xAxisStart, -0.6) -- (\xAxisEnd, -0.6);
\node[anchor=west] at (\xAxisEnd+0.2,-0.6) {Timestep, $k$};

% Braces above axis (match legacy)
\draw[decorate,decoration={brace,amplitude=4pt}] (7.0,0.05) -- (10.4,0.05)
  node[braceLabel, yshift=6pt] {Chunk consumption steps};
\draw[decorate,decoration={brace,amplitude=4pt}] (10.6,0.05) -- (12.9,0.05)
  node[braceLabel, yshift=6pt] {Steps consumed during inference};

% Small tick marks inside the brace ranges
\foreach \x in {7.15,7.55,7.95,8.35,8.75,9.15,9.55,9.95,10.35} {
  \draw[line width=0.5pt] (\x,-0.42) -- (\x,-0.78);
}
\foreach \x in {10.75,11.15,11.55,11.95,12.35,12.75} {
  \draw[line width=0.5pt] (\x,-0.42) -- (\x,-0.78);
}
\node[font=\scriptsize] at (7.12,-0.25) {$0_0$};
\node[font=\scriptsize] at (8.85,-0.25) {$\cdots$};
\node[font=\scriptsize] at (10.32,-0.25) {$0_k$};

% ----------------------------
% Helpers (macros)
% ----------------------------
% Draw a segmented rounded rectangle chunk at (x,y) with n slots.
\newcommand{\Chunk}[4]{%
  % #1 label (text), #2 x, #3 y, #4 n slots
  \pgfmathsetmacro{\w}{#4*\slotW}%
  \draw[chunkBox] (#2,#3) rectangle ++(\w,\slotH);
  \foreach \i in {1,...,#4} {%
    \draw[line width=0.45pt] (#2+\i*\slotW,#3) -- (#2+\i*\slotW,#3+\slotH);
  }%
  \node[anchor=west] at (#2+\w+0.25,#3+0.5*\slotH) {#1};
}

% Put a centered \cdots inside a chunk for readability (like the legacy figure).
\newcommand{\ChunkDots}[3]{%
  % #1 x, #2 y, #3 n slots
  \pgfmathsetmacro{\w}{#3*\slotW}%
  \node at (#1+0.5*\w,#2+0.5*\slotH) {$\cdots$};
}

% Fill an overlap region (light blue) across slot indices [i0, i1) in a chunk.
\newcommand{\Overlap}[5]{%
  % #1 x, #2 y, #3 i0 (0-based), #4 i1 (0-based), #5 height (defaults to \slotH)
  \pgfmathsetmacro{\x0}{#1 + #3*\slotW}%
  \pgfmathsetmacro{\x1}{#1 + #4*\slotW}%
  \path[overlapFill] (\x0,#2) rectangle (\x1,#2+#5);
}

\newcommand{\Pop}[2]{%
  \node[anchor=south] at (#1,#2) {\scriptsize pops};
  \draw[popArrow] (#1,#2-0.05) -- (#1,#2-0.55);
}

% ----------------------------
% Main frames (right): schedule/chunks and merge
% ----------------------------
% Frame 1: First chunk arrives and is merged into the schedule (conceptual).
\Chunk{$A_0$, first incoming action chunk}{7.1}{-1.95}{9}
\ChunkDots{7.1}{-1.95}{9}
\Pop{7.2}{-2.25}

% Show consumption by shifting the schedule (conceptual).
\Chunk{}{7.55}{-3.05}{9}
\ChunkDots{7.55}{-3.05}{9}
\Pop{7.65}{-3.35}

% During inference, some steps consumed (longer horizontal spacing to mimic legacy).
\Chunk{}{11.2}{-4.55}{7}
\ChunkDots{11.2}{-4.55}{7}
\Pop{11.25}{-4.85}

% Frame 2: Overlap/merge illustration.
% Incoming chunk A_t overlaps with remaining scheduled actions.
\Chunk{$A_t$, incoming action chunk}{11.0}{-6.55}{8}
\Chunk{$A_0$, current schedule}{11.8}{-7.45}{8}
\ChunkDots{11.0}{-6.55}{8}
\ChunkDots{11.8}{-7.45}{8}

% Overlap region (light blue), showing aggregated timesteps.
\Overlap{11.0}{-6.55}{2}{6}{\slotH}
\Overlap{11.8}{-7.45}{0}{4}{\slotH}

\node[labelText] at (13.1,-8.5) {Overlapping action timesteps are aggregated\\(aggregated actions in light blue)};

% Frame 3: Post-merge schedule (aggregated region visible).
\Chunk{}{12.2}{-9.75}{9}
\ChunkDots{12.2}{-9.75}{9}
\Overlap{12.2}{-9.75}{1}{7}{\slotH}
\Pop{12.25}{-10.05}
\node[font=\scriptsize, anchor=west] at (12.2,-10.45) {Action schedule (ordered by timestep)};

% ----------------------------
% Small callout: must_go triggers only when schedule is empty
% ----------------------------
\node[labelText, anchor=west] at (7.1,-10.95) {\scriptsize must\_go is set only when schedule is empty\\[-1pt]\scriptsize (armed when draining any new chunk)};

\end{tikzpicture}
\end{document}

